\chapter{Bunions (Hallux Valgus)}

\section*{Definition}
A bunion (hallux valgus) is a bony deformity at the base of the great toe, characterized by lateral deviation of the big toe and medial prominence of the first metatarsal head. It affects approximately 23 percent of adults and is more common in women.

\section*{When to See a Doctor}
\begin{itemize}
\item Visible bump on the inner side of the foot at the base of the big toe
\item Pain over the bunion that worsens with walking or tight shoes
\item Redness, swelling, or warmth over the joint (bursitis)
\item Numbness or tingling in the big toe
\item Difficulty fitting into regular shoes
\item Corns or calluses on the bunion or overlapping second toe
\item Restricted range of motion at the great toe joint
\item Progressive worsening of the deformity
\end{itemize}

\section*{How It Develops}
Bunions develop when the great toe gradually angles toward the other toes and the first metatarsal shifts outward. The prominence can rub against footwear, causing pressure, irritation, inflammation, and pain. Foot structure and inherited factors are important; footwear can worsen symptoms but is not the sole cause.

\section*{Causes}
\begin{itemize}
\item Genetic predisposition (inherited foot structure)
\item Tight, narrow, or high-heeled footwear
\item Flat feet (pes planus) with excessive pronation
\item Ligamentous laxity or hypermobility
\item Rheumatoid arthritis
\item Neuromuscular disorders affecting foot muscle balance
\item Pes cavus (high-arched feet)
\item Occupational factors (prolonged standing)
\item Female sex (four times more common than in men)
\end{itemize}

\section*{Related Symptoms}
\begin{itemize}
\item Foot pain (metatarsalgia)
\item Corns and calluses
\item Hammer toe deformity
\item Flat feet (pes planus)
\item Arthritis of the great toe joint (hallux rigidus)
\item Bursitis
\item Sesamoiditis
\end{itemize}

\section*{Medical Evaluation}
\begin{itemize}
\item Clinical examination (inspect deformity, assess range of motion)
 \item Weight-bearing X-rays when the diagnosis is uncertain, symptoms are significant, or surgery is being considered
\item Assess for associated deformities (hammer toes, metatarsalgia, flat feet)
\item Evaluate vascular status and sensation of the foot
\end{itemize}

\section*{Treatment Options}
\begin{itemize}
\item Conservative: footwear modification, orthotics, padding, ice, NSAIDs
\item Physical therapy for strengthening intrinsic foot muscles
 \item A clinician may consider an injection for selected inflammatory pain, but it does not correct the deformity
 \item Surgical correction if pain and functional limitation persist despite appropriate non-surgical care; surgery is not based on appearance alone
\item Various surgical techniques (osteotomy, arthrodesis, exostectomy)
\end{itemize}
\section*{Self-Care and Home Management}
\begin{itemize}
\item Wear wide-toed shoes with low heels
\item Use bunion pads (gel or foam) to cushion the prominence
\item Apply ice to reduce acute inflammation after prolonged walking
\item Take over-the-counter pain relievers as needed
\item Perform toe stretches and strengthening exercises
\item Use orthotic inserts to support the arch and redistribute pressure
\item Avoid activities that exacerbate pain until inflammation subsides
\end{itemize}

\section*{References}
\begin{thebibliography}{99}
\bibitem{bunions:ref1} Nix S, Smith M, Vicenzino B. ``Prevalence of Hallux Valgus in the General Population: A Systematic Review and Meta-Analysis.'' J Foot Ankle Res. 2010;3:21.
\bibitem{bunions:ref2} Coughlin MJ, Jones CP. ``Hallux Valgus: Demographics, Etiology, and Radiographic Assessment.'' Foot Ankle Int. 2007;28(7):759--777.
\bibitem{bunions:ref3} Mann RA, Coughlin MJ. ``Hallux Valgus: Etiology, Anatomy, Treatment and Surgical Considerations.'' Instr Course Lect. 1981;30:130--145.
\bibitem{bunions:ref4} Perera AM, Mason L, Stephens MM. ``The Pathogenesis of Hallux Valvis.'' J Bone Joint Surg Am. 2011;93(17):1650--1661.
\bibitem{bunions:ref5} Torkki M, Malmivaara A, Seitsalo S, et al. ``Surgery vs Orthosis vs Watchful Waiting for Hallux Valgus.'' JAMA. 2001;285(19):2474--2480.
\end{thebibliography}
